\section{LZ 编码：不假设分布的通用压缩}
\label{sec:07-Information-Theory/04-Source-Coding/06}

本章到目前为止，每一节的第一句都默认了同一件事：概率表已经摆在桌上。Kraft 用它算码长，Huffman 用它建树，算术编码用它切区间，率失真用它定下界。可你手上那个日志目录、那个源码仓库、那个 CSV，没有谁附带一份概率表。

先扫一遍统计频率再建表当然可以，但这条路有两个漏洞。一是频率表本身要编码传给解码端，小文件里它可能比正文还大；二是更根本的——很多数据的规律根本不在单符号频率上。英文文本里 \texttt{the} 反复出现，代码里 \texttt{for (} 反复出现，日志里时间戳的格式每行重来一遍。这些模式横跨多个符号，任何单符号频率表都看不见它们。

LZ（Lempel--Ziv）系列的做法是绕开显式的概率估计：把数据里重复出现的子串，换成一个指回先前位置的指针。

\subsection{重复的部分不必写第二遍}

LZ77 维护一个滑动窗口，窗口里是已经编码过的历史。每到一个新位置，就在历史中找与当前位置开头最长的匹配；找到了就输出``往回多远、抄多长''，找不到就直接输出这个字符本身。

\begin{center}
\begin{tikzpicture}[x=1cm,y=1cm,font=\small,>=stealth]
  \fill[blue!12] (0,0) rectangle (3.2,0.6);
  \fill[blue!28] (5.6,0) rectangle (8.8,0.6);
  \foreach [count=\i from 0] \c in {a,b,r,a,c,a,d,a,b,r,a} {
    \draw[black!50] (0.8*\i,0) rectangle (0.8*\i+0.8,0.6);
    \node[font=\footnotesize] at (0.8*\i+0.4,0.3) {\texttt{\c}};
  }
  \node[left,font=\footnotesize] at (-0.15,0.3) {输入};
  \draw[->,black!60] (7.2,0.72) -- (7.2,1.4) -- (1.6,1.4) -- (1.6,0.72);
  \node[above,font=\footnotesize,black!70] at (4.4,1.44) {往回数 \(7\) 个位置，照抄 \(4\) 个字符};
  \foreach [count=\i from 0] \c in {a,b,r,a,c,a,d} {
    \draw[black!50] (0.8*\i,-1.5) rectangle (0.8*\i+0.8,-0.9);
    \node[font=\footnotesize] at (0.8*\i+0.4,-1.2) {\texttt{\c}};
  }
  \fill[blue!28] (5.6,-1.5) rectangle (8.8,-0.9);
  \draw[black!50] (5.6,-1.5) rectangle (8.8,-0.9);
  \node[font=\footnotesize] at (7.2,-1.2) {指针 \((7,4)\)};
  \node[left,font=\footnotesize] at (-0.15,-1.2) {输出};
\end{tikzpicture}

\emph{第二次出现的 \texttt{abra} 不必重写，一个指针就够了。整个过程中没有出现概率表，压缩完全来自数据自己的重复结构。}
\end{center}

压缩效果因此取决于窗口大小和重复模式的间距：重复得越近越长，省得越多。文本、代码和日志的重复通常发生在几千字节以内，所以几 KB 到几十 KB 的窗口就已经很有效；窗口开大能捕捉更远的重复，内存和匹配搜索的时间也随之上去。

LZ78 换了一种记忆方式，用一本显式的字典代替滑动窗口：每当遇到``已见前缀加一个新符号''，就把这个新组合收进字典，输出的是``字典索引加新符号''。它的改进版 LZW 让字典初始就装好全部单符号，每次输出当前最长已见子串的索引，同时把该子串再接一个符号的新组合入库。两边的字典按同一条规则生长，解码端不需要接收字典，边解边建就同步了。gzip 用的 DEFLATE 是 LZ77 加 Huffman，GIF 与早期 ZIP 用 LZW，它们是同一思想的不同工程取舍。

\subsection{不写下分布，也能逼近熵率}

LZ 家族真正值钱的性质叫通用性。对平稳遍历的有限字母表信源，经典 LZ 方案在相应定理的条件下渐近码率趋于熵率 \(\bar H\)，而算法从头到尾不需要知道信源是什么。同一段代码压英文、压中文、压基因序列，都会自动逼近各自的熵率。

这件事初看有点反常——上一节刚说过压缩率就是交叉熵，交叉熵总要有个 \(Q\)。LZ 的 \(Q\) 藏在字典里。随着数据变长，字典把越来越多的常见模式收了进去，指针分配给频繁模式的比特越来越短，这等价于一个隐式的概率模型在不断变准。它没有把分布写在纸上，但它确实在学分布，而且学习的开销被摊进了同一条码流。

限定词不能删。平稳、遍历、字母表有限、序列足够长，四条缺一不可。短序列上字典还没长起来，指针开销就已经压过了收益；非平稳的数据流上，字典学到的是过去那一段的规律，未必对下一段有用。

\subsection{把压缩率当模型分数：能测什么，测不了什么}

上一节那条``压缩率就是交叉熵''的等式，在这里给出一种诱人的评测方式：把模型拿去压一段没见过的文本，压出来越小说明模型越好。Hutter Prize 直接用这个规则给通用压缩排名，语言模型报告的每字节比特数也是同一个量的另一种写法。

这套评测有它扎实的一面。它是端到端的：模型的每一个缺陷——概率算偏、上下文用得不好、罕见词处理粗糙——都会如实体现在字节数上，没有中间指标可以粉饰。它也无法作弊，因为解码端必须真的能还原原文，含糊的``差不多''不被接受。

局限同样明确。它测的是分布拟合的质量，仅此而已。一个把测试语料背下来的模型压缩率无可挑剔，但它什么也不会。压缩率对分布的整体形状敏感，对少数关键判断不敏感——模型在千分之一的样本上给出灾难性错误，字节数几乎不动。而且计分依赖于分词、字节还是字符，不同粒度的数字不能直接比。把它当作``模型有没有抓住数据的统计结构''的度量是恰当的，当作``模型有多好用''的度量就越界了。

\subsection{两类冗余，两种工具}

实际压缩器很少只用一种手段。DEFLATE 的结构说明了分工：LZ77 先把重复子串换成指针，去掉的是结构冗余；输出的字面量、长度与距离仍然频率不均，再交给 Huffman 去掉统计冗余。两级串联，因为它们看见的是不同的东西——一个看序列里的重复，一个看符号出现的偏斜。

反过来说，两种工具都失效的数据就是压不动的。加密输出、随机数、已经压过一遍的文件，既没有可用的重复，符号分布也接近均匀；LZ 找不到匹配还要为每个字面量额外付指针格式的开销，结果往往比原文更长。压缩器输出变大不是故障，是信息论上的必然——一个能压缩所有输入的算法不可能存在，因为长度为 \(n\) 的串有 \(2^n\) 个，短于 \(n\) 的串合起来才 \(2^n-1\) 个。

顺便澄清一个常被当作数据属性的判断：``这份数据不可压缩''其实从来不是数据的性质，而是``目前这个模型看不出其中的结构''的说法。一段伪随机数序列在 gzip 眼里是纯噪声，在知道种子和算法的人眼里可以压到几十个字节。熵是相对于模型的量，换模型就换下界。

无损压缩的账到此结清：熵是底线，前缀码把码长和概率对上号，Huffman 与算术编码把底线贴到可以忽略的距离，AEP 解释了底线为何不可逾越，率失真给出松开精确还原之后的新底线，而 LZ 说明连分布未知都不是障碍。这一整套推理有一个共同的前提——信道是干净的，写下去的比特一定原样读出来。真实的存储介质和通信链路不提供这个保证。下一章\hyperref[ch:058]{``有噪信道编码：噪声中为什么仍能可靠通信''}反过来做同一件事：不是删掉冗余，而是精心地加回冗余，让噪声吃掉一部分之后仍然能完整恢复。
